\documentclass[11pt,a4paper]{article}

% --- Input & fonts ---
\usepackage{iftex}
\ifPDFTeX\usepackage[utf8]{inputenc}\fi
\usepackage[T1]{fontenc}

% --- Math & theorems ---
\usepackage{amsmath,amssymb,amsthm}
\DeclareMathOperator{\Var}{Var}
\newcommand{\E}{\mathbb{E}}
\newcommand{\Prob}{\mathbb{P}}
\newcommand{\R}{\mathbb{R}}
\newcommand{\N}{\mathbb{N}}

% Currency symbol (for \euro)
\usepackage{eurosym}

% Financial notation
\newcommand{\ND}{\mathrm{ND}} % Net Debt
\newcommand{\EBITDA}{\mathrm{EBITDA}}
\newcommand{\EBITDAR}{\mathrm{EBITDAR}}
\newcommand{\ICR}{\mathrm{ICR}}
\newcommand{\Lev}{\mathrm{Lev}}
\newcommand{\Ifin}{I^{\mathrm{fin}}}
\newcommand{\Ilease}{I^{\mathrm{lease}}}
\newcommand{\CFCF}{C_{\text{FCF}}}
\newcommand{\pimin}{\pi_{\min}}
\newcommand{\pimax}{\pi_{\max}}
\newcommand{\CPI}{\mathrm{CPI}}
\newcommand{\FCF}{\mathrm{FCF}}

% Theorem environments (plain for results, definition style for defs/assumptions)
\theoremstyle{plain}
\newtheorem{theorem}{Theorem}
\newtheorem{proposition}{Proposition}
\newtheorem{lemma}{Lemma}
\theoremstyle{definition}
\newtheorem{definition}{Definition}
\newtheorem{assumption}{Assumption}
\renewcommand\theassumption{A\arabic{assumption}}

% --- Figures, tables, algorithms, typography ---
\usepackage{graphicx}
% Prefer local paper figures first, then shared analysis figures
\graphicspath{{paper/figures/v2/}{figures/v2/}}
\DeclareGraphicsExtensions{.pdf,.png,.jpg,.jpeg}
\usepackage{subcaption}
\usepackage{booktabs}
\usepackage{microtype}
\usepackage{float}                 % [H] exact placement
\usepackage[section]{placeins}     % keep floats within sections

% Allow pasted Unicode math symbols
\ifdefined\DeclareUnicodeCharacter
\DeclareUnicodeCharacter{2265}{$\ge$} % ≥
\DeclareUnicodeCharacter{2264}{$\le$} % ≤
\DeclareUnicodeCharacter{2212}{$-$}   % −

\fi

% --- Bibliography (author–year) ---
\usepackage[authoryear,round]{natbib}
\setlength{\bibsep}{0pt plus 0.3ex}

% --- Page & links ---
\usepackage{geometry}
\geometry{margin=1in}

\usepackage{hyperref}
\hypersetup{
  colorlinks=true,
  linkcolor=blue,
  filecolor=magenta,
  urlcolor=cyan,
  citecolor=red,
  pdftitle={Covenant Screening Under IFRS-16: Assumption-Bounded Analytic Approximations and Simulation-Based Evaluation},
  pdfauthor={Aniket Bhardwaj},
  pdfkeywords={leveraged buyout, IFRS-16, covenant screening, synthetic evaluation, computational finance},
  pdfcreator={LaTeX}
}
\urlstyle{same}

% --- Clever references: load after hyperref (with fallback if missing) ---
\IfFileExists{cleveref.sty}{%
  \usepackage[capitalise,nameinlink]{cleveref}
  \crefname{proposition}{Proposition}{Propositions}
  \crefname{theorem}{Theorem}{Theorems}
  \crefname{lemma}{Lemma}{Lemmas}
  \crefname{definition}{Definition}{Definitions}
  \crefname{assumption}{Assumption}{Assumptions}
  \crefname{figure}{Figure}{Figures}
  \crefname{table}{Table}{Tables}
  \crefname{section}{Section}{Sections}
}{%
  % Minimal fallback to avoid compile failure when cleveref is unavailable
  \newcommand{\cref}[1]{\ref{#1}}
  \newcommand{\Cref}[1]{\ref{#1}}
}

% Title and authors
\title{Covenant Screening Under IFRS-16: \\
Assumption-Bounded Analytic Approximations and Simulation-Based Evaluation}
\author{
Aniket Bhardwaj\\
\textit{Independent Researcher}\\
\texttt{bhardwaj.aniket2002@gmail.com}
}

\usepackage{tabularx}
\usepackage{xurl}
\IfFileExists{paper/generated/metrics_table.tex}{\def\generatedpath{paper/generated/}}{\def\generatedpath{generated/}}
\makeatletter
\newcommand{\paperinput}[1]{\@@input\generatedpath#1 }
\makeatother
\date{Revised manuscript}
\begin{document}
\maketitle
\begin{abstract}
We compare IFRS-16-inclusive and stipulated frozen-GAAP covenant calculations
using a reduced-form analytic screening model and a full annual financing
simulation. The latter funds opening cash through closing sources and uses and
reconciles mandatory principal, revolver drawings, cash sweeps and payment
default. Deterministic screening statements apply only to matched reduced-form
paths with established error budgets and positive denominator floors; empirical
cross-model discrepancies also reflect structural, timing and input differences.
In the current seed-42 synthetic benchmark of 200 five-year scenarios, the
analytic ranking score has AUC 0.9472 (scenario-bootstrap 95\% interval
[0.8974, 0.9828]), but the 0.5 classification threshold has an 85\% false-negative
rate and a zero false-positive rate. Twenty scenarios financially fail, including
17 payment defaults. Mean scenario RMSE is 0.8084 for leverage and 2.3157 for
coverage; minimum-headroom RMSE is 1.4884. A separate Accor accounting-convention
illustration reproduces supplied reported/reconstructed inputs and hypothetical
thresholds, without claiming predictive validation. The score is not a calibrated
default probability, and the benchmark supplies no equity-return or optimization
uplift result. Code, input assumptions, numerical artifacts and six regenerated
figures accompany the revised manuscript.
\end{abstract}

\textbf{Keywords:} IFRS-16, leveraged buyout, covenant screening, synthetic benchmark\\
\textbf{JEL Classification:} G24, G32, C61

\section{Introduction}
Lease capitalization changes the quantities entering debt covenants, while a credit
agreement's definitions determine which accounting effects enter its tests. IFRS~16,
issued in 2016 and effective for annual periods beginning in 2019, brings most lessee
lease commitments onto the balance sheet, subject to its exemptions
\citep{ifrs2016leases}. Comparing a lease-inclusive covenant calculation with a
stipulated rent-adjusted alternative is therefore useful for lease-intensive
transactions. Neither calculation by itself establishes the terms of an actual loan.

This paper studies assumption-bounded analytic screening alongside an explicit
annual financing simulation. These are two materially different model families.
The analytic model provides transparent debt and coverage paths. The simulation
reconciles closing funding, operating cash, debt service, revolver capacity and
cash retention. Differences between their outputs combine approximation error,
structural-model differences, timing differences and input/convention differences.
An empirical comparison cannot isolate the first component without a separate,
matched-model experiment.

\subsection{Contributions and scope}
The contribution is a reproducible characterization of this screening workflow,
conditional mathematical statements about matched reduced-form paths, and an
evaluation on synthetic scenarios. We distinguish a ranking score from an
operational decision rule, and distinguish covenant compliance from the broader
financial-failure union. An Accor SA accounting-convention illustration retains
the link to a named company while making reconstructed inputs and hypothetical
thresholds explicit. The present evaluation does not fit a borrower prediction
model or estimate the return from an optimized covenant package.

The revision replaces historical empirical claims after financing corrections.
In particular, retained opening cash is a funded use, so sponsor equity and any
return calculated from it must use the repaired closing bridge. Historical IRR
tables cannot be carried forward. The title accordingly emphasizes screening and
simulation-based evaluation rather than an empirically demonstrated optimization
or a universal deterministic guarantee.

\section{Related Work and Background}
\subsection{Covenants and control rights}
\citet{dichev2002quality} investigate the debt covenant hypothesis using a large
sample of private lending agreements. \citet{nini2009creditor} study creditor
control rights and investment restrictions outside payment default. These studies
motivate careful separation of covenant conditions from missed debt service; they
do not validate the synthetic data-generating process used here. The present
exercise concerns model behavior under stipulated inputs, rather than estimating
contracting effects in those borrower populations.

\subsection{Lease reporting and covenant conventions}
IFRS~16 specifies lease reporting, whereas contractual covenant definitions may
adjust accounting quantities. We use the repository's two named conventions as
explicit formulas. Its frozen-GAAP helper adds supplied rent to supplied EBITDA
for coverage and excludes leases from net debt. This is a stipulated rent-adjusted
coverage convention, not an automatic reconstruction of historical-GAAP accounts.
Historical operating-lease EBITDA deducts operating rent; IFRS-16 EBITDA generally
does not include that same operating rent charge. Inputs must therefore be aligned
with the intended credit agreement before using either helper.

\section{Notation and Constants}
\begin{table}[H]
\centering\small
\caption{Notation and units. Superscript $a$ identifies the analytic family.}
\begin{tabularx}{\textwidth}{l l X}
\toprule Symbol & Unit & Definition\\\midrule
$t,T$ & years & Annual index and horizon; empirical $T=5$.\\
$B_t,R_t,L_t,C_t$ & currency & Closing term debt, revolver, lease liability and cash.\\
$E_t,V_t$ & currency/year & EBITDA and revenue.\\
$I_t^d,I_t^L$ & currency/year & Financial interest and lease interest.\\
$A_t,U_t,S_t$ & currency/year & Actual mandatory amortisation, unpaid amount and sweep.\\
$X_t^a,X_t^\ell,Q_t$ & currency/year & Amortisation draw, liquidity draw and revolver repayment.\\
$s,q,\eta$ & fractions & Sweep, lease principal and lease-addition coefficients.\\
$c^{\rm lev},c^{\rm icr}$ & ratios & Maximum leverage and minimum coverage.\\
$h_t,H$ & ratio points & Minimum of two covenant margins; minimum over time.\\
$\epsilon_Z$ & units of $Z$ & Assumed or established error budget for component $Z$.\\
\bottomrule\end{tabularx}
\end{table}
Combining leverage and coverage margins by their numerical minimum is the
implemented screening convention. The two margins have different economic
interpretations despite both being dimensionless ratios; this minimum is not a
monetary distance to default.

\section{Model Framework}
\subsection{Dual covenant conventions}
For positive EBITDA and positive applicable interest denominators, the project
calculates
\begin{align}
\ND_t^{16}&=B_t+R_t+L_t-C_t,&
\Lev_t^{16}&=\frac{\ND_t^{16}}{E_t},&
\ICR_t^{16}&=\frac{E_t}{I_t^d+I_t^L},\label{eq:ifrs}\\
\ND_t^{F}&=B_t+R_t-C_t,&
\Lev_t^{F}&=\frac{\ND_t^{F}}{E_t},&
\ICR_t^{F}&=\frac{E_t+P_t}{I_t^d},\label{eq:frozen}
\end{align}
where $P_t$ is supplied rent in the frozen-GAAP helper. The full synthetic benchmark
evaluates only the IFRS-16-style view in \eqref{eq:ifrs}; the separate Accor
illustration evaluates both. Nonpositive EBITDA is a failure condition rather than
evidence of favorable negative leverage. The Accor helpers return undefined ratios
in that state. The benchmark instead checks the failure flag separately and floors
the interest denominator at $10^{-6}$ in calculating simulated coverage.

\subsection{Reduced-form analytic screening}
The implemented paths use $E_t^a=E_0^a(1+g)^t$ and
$F_t^a=(\alpha-\kappa)E_t^a$, with weighted financial rate
$r_d=w r_{\rm sen}+(1-w)r_{\rm mezz}$. They obey
\begin{align}
D_{t+1}^a&=\max\{0,(1+r_d)D_t^a-sF_t^a\},\label{eq:analyticdebt}\\
L_{t+1}^a&=\max\{0,(1-q)L_t^a+\eta E_t^a\},\label{eq:analyticleases}\\
C_0^a&=C_{\rm initial}^a,\qquad C_t^a=C_{\min}^a\quad(t\ge1).
\end{align}
Equation~\eqref{eq:analyticleases} is the run-off mode used by the benchmark.
The alternative steady-state mode sets $L_t^a=\lambda E_t^a$. Where the debt floor
does not bind, the familiar linear-recurrence expression is
\begin{equation}
D_t^a=D_0^a(1+r_d)^t-s\sum_{k=0}^{t-1}(1+r_d)^{t-1-k}F_k^a.
\label{eq:closedform}
\end{equation}
After a floor binds, the implemented recursion, rather than an unrestricted
closed-form extrapolation, defines the path. Analytic coverage uses closing debt
and lease balances. Interest capitalization in \eqref{eq:analyticdebt} and fixed
cash are reduced-form assumptions; there is no reconciled mandatory-amortisation
or revolver waterfall in this family.

\subsection{Full simulation: closing and operating cash}
For stipulated debt-free purchase use $V_{\rm purchase}$, fee fraction $f$ and
retained opening cash $C_0$, sponsor funding is
\begin{equation}
E_{\rm sponsor}=V_{\rm purchase}(1+f)+C_0-B_0.\label{eq:closing}
\end{equation}
Closing debt plus sponsor equity equals total uses. Opening cash is not a source
of acquisition funding. Debt above uses is rejected; zero sponsor equity remains
exactly zero, with undefined IRR and MOIC. Default assumptions provide a mechanical
check: purchase 1,000, fees 30 and cash 40 are funded by debt 450 and equity 620.
These are model currency units, not observed transaction amounts.

Revenue in year 1 equals the supplied starting revenue; growth begins in year 2:
$V_t=V_0(1+g)^{t-1}$ and $E_t=mV_t$. Using opening balances,
\begin{align}
I_t^d&=r_d^{s}(B_{t-1}+R_{t-1}),& I_t^L&=r_L^{s}L_{t-1},\\
\mathrm{Tax}_t&=\max\{0,\tau(E_t-dV_t-I_t^d-I_t^L)\},\\
O_t&=E_t-\Delta WC_t-\mathrm{Tax}_t-I_t^d-I_t^L-kV_t,\\
K_t&=C_{t-1}+O_t-qL_{t-1},&
L_t&=\max\{0,(1-q)L_{t-1}+\eta V_t\}.
\end{align}
Here $K_t$ is cash before financing. Working capital is a fixed revenue fraction;
the initial working-capital balance makes year-1 $\Delta WC$ zero. Lease additions
are noncash and proportional to current revenue. This stylized schedule is not a
lease-by-lease IFRS valuation or a complete financial-statement model.

\subsection{Full simulation: financing waterfall}
Let $M_t=\min\{M,B_{t-1}\}$ denote scheduled principal and
$U_t^R=\max\{0,R_{\max}-R_{t-1}\}$ available revolver capacity. The waterfall is
\begin{align}
A_t^{\rm cash}&=\min\{M_t,\max(0,K_t)\},\\
X_t^a&=\min\{M_t-A_t^{\rm cash},U_t^R\},& A_t&=A_t^{\rm cash}+X_t^a,\\
U_t&=M_t-A_t,& K_t'&=K_t-A_t^{\rm cash},\\
X_t^\ell&=\min\{\max(0,C_{\min}-K_t'),\max(0,U_t^R-X_t^a)\}.
\end{align}
Thus amortisation drawings go directly to the term lender. They do not increase
retained cash. A negative operating balance remains negative unless a separate
liquidity draw funds it. After these draws, excess cash repays the revolver first.
If $K_t''=K_t'+X_t^\ell-Q_t$ is cash after that repayment, the optional sweep is
\begin{equation}
S_t=\min\{\max(0,K_t''-C_{\min}),s\max(0,K_t),\max(0,B_{t-1}-A_t)\}.
\end{equation}
The resulting accounting identities are
\begin{align}
C_t&=C_{t-1}+O_t-qL_{t-1}-A_t-S_t+X_t^a+X_t^\ell-Q_t,\\
B_t+R_t&=B_{t-1}+R_{t-1}-A_t-S_t+X_t^a+X_t^\ell-Q_t.
\end{align}
Unpaid mandatory principal stays outstanding. Payment default means $U_t>0$.
Reserve failure means $C_t<C_{\min}$, not a legal insolvency finding. Mandatory
principal takes priority over the reserve. Simulation continues after failure for
diagnostics, without a recovery or arrears catch-up model.

\subsection{Covenant design and an unevaluated return objective}
The historical conceptual objective can be written, for a decision vector $x$ and
explicit uncertainty distribution $\mathcal P$, as
\begin{align}
\sup_x\quad &\E_{\theta\sim\mathcal P}[\mathrm{IRR}(x;\theta)]\label{eq:objective}\\
\text{subject to}\quad &\Lev_t^a(x)+\epsilon_{\rm lev}(t)\le c^{\rm lev},\nonumber\\
&\ICR_t^a(x)-\epsilon_{\rm icr}(t)\ge c^{\rm icr},\quad t=1,\ldots,T.\nonumber
\end{align}
This is a possible reduced-form design formulation, not a solved empirical
optimization in the current benchmark. If only thresholds change and financing
and cash flows are fixed, an individual deal's return does not change. Claims of
return improvement require a defined financing decision and reoptimized feasible
set. The implementation's sponsor vector is
$[-E_{\rm sponsor},0,\ldots,0,E_{\rm exit}]$, with exit EBITDA multiple less sale
costs, financial debt and leases, plus cash. No equity-return output is emitted
by the benchmark, so no expected-IRR or IRR-improvement estimate is reported here.

\subsection{Bayesian hierarchical calibration: separate research framework}
The earlier framework contemplated bounded-support hierarchical parameters.
For bounds $a_j<b_j$, one possible specification is
\begin{align}
z_{ij}&\sim\mathcal N(\mu_j,\sigma_j^2),&
\theta_{ij}&=a_j+(b_j-a_j)\{1+\exp(-z_{ij})\}^{-1},\\
\log\Lambda_i&\sim\mathcal N(\mu_\Lambda,\sigma_\Lambda^2),&
\sigma_j&\sim\mathrm{HalfNormal}(s_j).
\end{align}
Growth shocks and cross-parameter dependence would require additional specified
distributions. These equations preserve the conceptual distinction between
parameter uncertainty and financial dynamics. They are not a description of a
posterior estimated for the present results. The separate calibration module has
optional PyMC and MAP/Laplace paths; the benchmark calls neither. It samples
clipped scenario parameters directly. Consequently this paper reports no
posterior credible frontier or empirical benefit from Bayesian calibration.

\section{Deterministic Approximation Bounds}
\subsection{Scope of the guarantees}
We retain deterministic reasoning only for a pair of matched reduced-form paths
with established component error budgets and positive denominators. A reference
path below is not automatically the full financing simulation. Assumptions in
Appendix~\ref{app:assumptions} are requirements to be checked, not consequences of
the current synthetic sampler. In particular, no universal financial-default
guarantee follows from bounding two covenant ratios.

\begin{proposition}[Conditional ratio screening]\label{prop:screening}
Suppose reference and analytic EBITDA and interest satisfy
$E,E^a\ge e>0$ and $I,I^a\ge i>0$ at a given date. Suppose
$|E^a-E|\le\epsilon_E$, $|N^a-N|\le\epsilon_N$,
$|I^a-I|\le\epsilon_I$, and known upper bounds are
$|E^a|\le\overline E$ and $|N^a|\le\overline N$. Then
\begin{align}
|\ICR^a-\ICR|&\le\epsilon_{\rm icr}
 :=\frac{\epsilon_E}{i}+\frac{\overline E\epsilon_I}{i^2},\\
|\Lev^a-\Lev|&\le\epsilon_{\rm lev}
 :=\frac{\epsilon_N}{e}+\frac{\overline N\epsilon_E}{e^2}.
\end{align}
For matched debt and lease definitions and fixed rates,
$\epsilon_N=\epsilon_D+\epsilon_L+\epsilon_C$ and
$\epsilon_I=r_d\epsilon_D+r_L\epsilon_L$ suffice.
\end{proposition}
Cash cannot be omitted from net-debt error unless it is identical between paths.
Different interest timing or rates require their own additional error budgets.
For the rent-adjusted coverage numerator, a rent error must likewise enter the
numerator budget. The proposition is deliberately conditional; the constants
emitted by \texttt{analytic\_bounds.py} do not establish its premises.

\begin{proposition}[Feasible-set monotonicity]\label{prop:frontier}
Fix the cash-flow mapping, testing dates, coverage threshold and objective. If
$c_1^{\rm lev}\le c_2^{\rm lev}$, the feasible sets defined by the same screening
inequalities satisfy $\mathcal F(c_1^{\rm lev})\subseteq\mathcal F(c_2^{\rm lev})$.
Their supremum objective values are therefore non-decreasing, when defined.
\end{proposition}
This is set inclusion, not a claim that relaxing a covenant changes a fixed
transaction's IRR or induces a specific borrower response.

\begin{proposition}[Conditional covenant certification]\label{prop:certification}
Under Proposition~\ref{prop:screening} at every testing date, let
$h_t^a=\min\{\ICR_t^a-c^{\rm icr},c^{\rm lev}-\Lev_t^a\}$ and
$\epsilon_{\max}=\max_t\max\{\epsilon_{\rm icr}(t),\epsilon_{\rm lev}(t)\}$.
If $h_t^a>\epsilon_{\max}$ at every date, the reference path strictly satisfies
both ratio covenants at those dates.
\end{proposition}
This statement does not certify reserve adequacy, payment of principal, or the
outcomes of the unmatched full simulation. No empirical certification percentage
is reported.

\section{IFRS-16 LBO Benchmark Generator}
\subsection{Synthetic population and executable protocol}
The authoritative procedure is \path{analysis/run_benchmark.py} in the
repository \citep{repository}. With seed 42,
it samples 200 scenarios from five synthetic operator templates with replacement,
each with five annual observations. The 20-scenario smoke mode is only a separate
software check. The metadata file \texttt{scenario\_parameters.csv} is checksummed
but does not control sampling. There is no audited borrower or transaction panel.

For each scenario, the sampler draws a base, downside or distressed regime with
probabilities $(0.5,0.3,0.2)$. Growth is initially normal around the operator's mean
with its supplied standard deviation and clipped to $[-0.12,0.18]$; margin has
standard deviation 0.02 and is clipped to $[0.10,0.35]$. Sweep has standard
deviation 0.04 and range $[0.30,0.85]$. Lease principal and additions have standard
deviations 0.01 and 0.003 and ranges $[0.05,0.20]$ and $[0,0.05]$ respectively.
Downside then subtracts 0.04 from growth and 0.03 from margin, multiplies debt by
1.05, leases by 1.04 and cash by 0.85. Distress uses respective changes
$-0.08,-0.06,1.20,1.10,0.60$. Final growth and margin floors are $-0.12$ and 0.08.
Draws are constant scenario parameters, not five new annual random shocks.

\begin{table}[H]\centering\small
\caption{Material distinctions in the actual benchmark model pair.}
\begin{tabularx}{\textwidth}{l>{\raggedright\arraybackslash}X>{\raggedright\arraybackslash}X}\toprule
Feature & Analytic screening & Full simulation\\\midrule
Initial EBITDA & Operator template EBITDA & Revenue times drawn margin\\
Growth timing & Applied at time 1 & Starts in year 2\\
Cash & Initial 25; fixed 10 thereafter & Regime-adjusted template cash; reserve 25\\
Financial interest & Weighted 0.062; closing debt & 0.06; opening term debt and revolver\\
Lease interest & 0.045; closing leases & 0.05; opening leases\\
Lease additions & Prior EBITDA basis & Current revenue basis\\
Debt service & Capitalization and FCF sweep & Cash interest; mandatory principal 30; revolver\\
Revolver capacity & None & 100 distressed; 200 otherwise\\
Operating deductions & $(\alpha-\kappa)E^a$ & Explicit tax, working capital and capex\\
\bottomrule\end{tabularx}\label{tab:families}
\end{table}
Analytic $\alpha$ is 0.72 in base and 0.68 otherwise; $\kappa$ is 0.10 in distress
and 0.08 otherwise. Full-simulation depreciation, tax, capex and working-capital
fractions are 0.03, 0.25, 0.04 and 0.02. Full scenario assumption records are
archived alongside paths, so the table is not a substitute for executable inputs.

\subsection{Failure outcome, score and uncertainty}
The outcome is the union, over all five years, of nonpositive EBITDA, payment
default, reserve funding deficit, leverage above 6.0 or coverage below 1.8.
Equality at a covenant threshold is not a simulated breach. The minimum analytic
headroom and ranking score are
\begin{equation}
H^a=\min\{6-\max_t\Lev_t^a,\min_t\ICR_t^a-1.8\},\qquad
z=(1+\exp(2H^a))^{-1}.\label{eq:score}
\end{equation}
The score is not a calibrated probability of default. Classification declares
failure for $z\ge0.5$, including exactly zero analytic headroom. Thus the decision
boundary's equality convention differs from the simulation's strict breach test.
No model fitting, train/test split, cross-validation or threshold optimization
is used in this evaluation.

AUC is computed across scenarios. Its percentile 95\% interval uses 500 scenario
bootstrap attempts with seed 42, skipping single-class resamples. This is not an
operator-clustered interval. Repeated scenarios share five template assumptions,
so the interval does not establish out-of-sample borrower performance. There are
no RMSE confidence intervals, permutation tests, calibration curves or Brier
scores in the current report. Execution exceptions and a single-class full sample
abort reporting instead of silently removing scenarios.

\section{Empirical Validation}
\subsection{Analytic versus simulation diagnostics}
\begin{figure}[H]\centering
\includegraphics[width=\textwidth]{F12_theoretical_guarantees}
\caption{Replacement for the historical theoretical-guarantees plot: signed
analytic-minus-simulation leverage and ICR differences over 1,000 scenario-years.
Zero is an equality reference, not an error bound. The distributions combine
structural, timing, input and approximation differences; they do not test a
universal deterministic guarantee. Source: archived regenerated scenario paths,
seed 42; exact source revision in Appendix~\ref{app:computation}.}
\label{fig:theoretical_bounds}\end{figure}

\begin{table}[H]\centering
\caption{Current cross-model discrepancy metrics, in ratio points.}
\begin{tabular}{lr}\toprule Metric & Value\\\midrule
\paperinput{metrics_table.tex}
\bottomrule\end{tabular}\label{tab:errors}\end{table}
Leverage and ICR MAE are averaged across scenarios after computing each scenario's
five-year MAE. Their reported RMSE is the mean of the 200 scenario-level
five-year RMSEs, not the pooled root mean square across all 1,000 observations.
Minimum headroom is one scalar per scenario; its RMSE is the square root of the
mean of 200 squared headroom differences. This distinction is necessary for exact
replication. The old headroom table used a different and unreproduced protocol;
Table~\ref{tab:errors} is a new evaluation, not an update to that table's cells.

\subsection{Benchmark performance analysis}
The AUC is 0.947222 with a scenario-bootstrap 95\% interval
$[0.897410,0.982827]$. At the unchanged 0.5 threshold the false-positive rate is
0 and the false-negative rate is 0.85. There are 3 true positives, 17 false
negatives, 180 true negatives and no false positives. Good ordering of failures
relative to nonfailures therefore coexists with poor failure detection at this
particular threshold. It would be inappropriate to describe the classifier as
conservative certification of the full financing model.

\begin{figure}[H]\centering
\includegraphics[width=\textwidth]{F14_method_comparison}
\caption{Current method comparison: the analytic score's ROC against full-simulation
financial-failure labels, and cross-model MAE/RMSE. The diagonal is a chance
reference, not a reproduced traditional or Bayesian comparator. Error aggregation
follows Table~\ref{tab:errors}; headroom and individual ratio errors measure
different quantities. Source: benchmark JSON, seed 42.}
\label{fig:benchmark_performance}\end{figure}

\begin{figure}[H]\centering
\includegraphics[width=\textwidth]{F13_benchmark_overview}
\caption{Synthetic benchmark overview: realized regime counts and operator-template
sampling frequencies among 200 scenarios. These are synthetic archetypes, not a
geographic borrower dataset. Source: benchmark scenario records, seed 42.}
\label{fig:benchmark_overview}\end{figure}

\subsection{Sensitivity and certification scope}
The benchmark fixes annual testing and the two stated covenant thresholds. It
does not evaluate quarterly testing, equity cures, hedging interventions,
interest-rate floors or optimized covenant packages. The historical IRR
sensitivity table has no reproducible counterpart and is not reported as a
current result. A future sensitivity experiment must define interventions,
matched financing inputs and a sponsor-return protocol before computing effects.
Likewise, heuristic diagnostic envelopes in the JSON are not established
component budgets for Proposition~\ref{prop:screening}. No bound utilization,
false non-certification, or time-point certification rate is inferred from them.

\subsection{Breach composition analysis}
Twenty of 200 scenarios financially fail; 180 do not. The mutually exclusive
priority labels identify 17 payment defaults and 3 other combined reserve/covenant
failures. A payment-default label may also involve reserve and covenant failures.
The separate flags in Figure~\ref{fig:breach_composition} preserve that overlap.
The JSON field \texttt{failed\_scenario\_count} is zero: it counts execution
exceptions, not these 20 financial failures.
\begin{figure}[H]\centering
\includegraphics[width=\textwidth]{F16_breach_composition}
\caption{Financial-failure composition. Left: disjoint priority labels, with payment
default preceding combined reserve/covenant failure. Right: overlapping payment,
reserve and covenant flags; these counts must not be added as if disjoint.
Source: full benchmark JSON, seed 42. No base-versus-optimized intervention is
claimed.}
\label{fig:breach_composition}\end{figure}

\subsection{Failure-mode diagnostics}
\begin{figure}[H]\centering
\includegraphics[width=\textwidth]{F15_failure_modes}
\caption{Analytic ranking-score distributions by realized financial failure and
the confusion matrix at 0.5. Actual simulation outcomes index rows; analytic
classifications index columns. The 17 missed failures make the distinction
between ranking discrimination and threshold performance visible. Source:
benchmark scenario records, seed 42.}
\label{fig:failure_modes}\end{figure}
The evidence describes failure detection conditional on the current DGP. It does
not identify a causal failure mechanism from growth volatility or validate old
percentages assigned to assumption-edge violations. Both model structure and
starting inputs can explain discrepancies. Attribution requires controlled
one-change-at-a-time or matched-model experiments beyond this evaluation.

\section{Accor SA Dual-Convention Analysis}\label{sec:accor}
\subsection{Illustrative real-company inputs}
We retain an illustrative real-company accounting-convention case study using
the repository's reported/reconstructed Accor inputs for 2019--2023. This is not
predictive validation. The CSV includes report, page and reconstruction notes;
these annotations are provenance claims in the repository, not an independent
audit of the cited filings. We reproduce the supplied values without asserting
that each value is an audited consolidated-statement line item. The old
2018--2022 table is not the input to the current script.

\begin{table}[H]\centering\small
\caption{Stipulated Accor case-study inputs (EUR millions) from
\texttt{data/case\_study/accor.csv}. Rent and interest are supplied inputs.}
\begin{tabular}{rrrrrrr}\toprule
Year & Revenue & EBITDA & Net debt & Leases & Rent & Interest\\\midrule
\paperinput{accor_inputs.tex}
\bottomrule\end{tabular}\label{tab:accor_data}\end{table}
The case-study script allocates supplied net debt 80/20 to two financial-debt
fields, sets cash to zero because the supplied quantity is already net, and sets
the financial rate to interest expense divided by net debt. The split therefore
does not change aggregate financial interest. It assumes a lease rate of 4\%
and uses supplied lease expense as rent. Including leases on top of the supplied
net debt is a stipulated bridge; it requires that the supplied net-debt field
exclude those leases for an economic interpretation without double counting.
The script does not independently establish that reporting convention from filings.

\subsection{Dual-convention ratios and hypothetical thresholds}
\begin{table}[H]\centering\small
\caption{Regenerated Accor ratios. F denotes the stipulated frozen-GAAP helper;
16 denotes the lease-inclusive helper. Undefined 2020 ratios are not zeros.}
\begin{tabular}{rrrrr}\toprule
Year & ICR F & ICR 16 & Leverage F & Leverage 16\\\midrule
\paperinput{accor_ratios.tex}
\bottomrule\end{tabular}\label{tab:accor_ratios}\end{table}
Across the four positive-EBITDA years, the arithmetic mean ICR changes from
4.764x in the stipulated frozen-GAAP calculation to 2.464x in the IFRS-16
calculation. Mean leverage changes from 4.559x to 9.323x. The negative-EBITDA
2020 observation is excluded from these ratio means and explicitly classified as
a failure in both views. These quantities describe the supplied inputs and
stipulated conventions; they are not estimates of a predictive model's accuracy.

\begin{figure}[H]\centering
\includegraphics[width=\textwidth]{accor_case_study}
\caption{Regenerated Accor accounting-convention illustration. Dashed thresholds
(ICR at least 4.0x and leverage at most 3.5x) are hypothetical, not Accor's actual
contractual covenants. The shaded 2020 year has negative EBITDA and undefined
ratios; gaps are intentional. Source: repository case-study CSV and the current
case-study helper, with its stipulated 4\% lease rate.}
\label{fig:accor_ratios}\end{figure}
Including the 2020 failure, the hypothetical tests fail in 5/5 years under the
lease-inclusive calculation and 3/5 under the stipulated frozen-GAAP view.
The example demonstrates sensitivity to definitions. It does not establish
actual contractual breaches, prove that one convention protects a borrower from
distress, or show that a redesigned covenant would improve investment returns.

\section{Limitations and Future Work}
The benchmark is synthetic, uses five archetypes and contains no audited
transaction dataset. Its AUC and interval are conditional on the specified DGP,
seed, outcome definition and model assumptions. There is no calibrated
probability-of-default interpretation or independent borrower test set.
The bootstrap samples scenarios rather than operator clusters and does not
quantify uncertainty about selecting the templates or the financial model.

Structural, timing and input differences between the two model families preclude
interpreting observed discrepancies as pure analytic approximation error.
The conditional mathematical guarantees require positive denominator floors and
component budgets that the benchmark does not establish. Cash, payment default
and reserve adequacy require explicit treatment beyond ratio certification.

The full model has simplified taxes and lease schedules, no complete loss
carryforward, no separate right-of-use depreciation schedule, no default recovery
model and no arrears catch-up. Negative cash and negative exit residuals are
diagnostic; the latter are not limited-liability equity payoffs. Comprehensive
input-domain validation is not implemented. The Accor exercise is illustrative
rather than predictive validation, and reconstructed inputs require independent
filing reconciliation before making company-specific economic claims.

Future work should first match the model inputs, dates, rates and cash
definitions, then isolate approximation components. A separate empirical study
could evaluate operator-held-out or transaction-held-out performance and score
calibration. Return optimization would require an explicit financing decision,
default/recovery treatment and a reproducible return protocol. These extensions
are research tasks, not findings of this revision.

\section{Conclusion}
The repaired engine supports a transparent synthetic evaluation of analytic
covenant screening against explicit financing dynamics. Ranking discrimination
and threshold performance give materially different assessments of this
workflow. Conditional mathematical screening and the Accor convention example
remain useful when their assumptions and input provenance are explicit. The
present evidence does not demonstrate calibrated borrower default probabilities,
universal simulation error bounds or an equity-return uplift.

\section*{Acknowledgments}
The implementation and reproducibility artifacts use the open-source scientific
Python and LaTeX ecosystems.

\begin{thebibliography}{99}
\bibitem[Dichev and Skinner(2002)]{dichev2002quality}
Dichev, I.~D., and Skinner, D.~J. (2002).
Large-sample evidence on the debt covenant hypothesis.
\emph{Journal of Accounting Research}, 40(4), 1091--1123.
\url{https://doi.org/10.1111/1475-679X.00083}.
\bibitem[IFRS Foundation(2016)]{ifrs2016leases}
IFRS Foundation (2016). \emph{IFRS 16: Leases.}
International Accounting Standards Board.
\url{https://www.ifrs.org/issued-standards/list-of-standards/ifrs-16-leases/}.
\bibitem[Nini et~al.(2009)]{nini2009creditor}
Nini, G., Smith, D.~C., and Sufi, A. (2009).
Creditor control rights and firm investment policy.
\emph{Journal of Financial Economics}, 92(3), 400--420.
\url{https://doi.org/10.1016/j.jfineco.2008.04.008}.
\bibitem[Bhardwaj(2026)]{repository}
Bhardwaj, A. (2026). \emph{IFRS-16 LBO Engine}, source code and
reported/reconstructed Accor case-study inputs.
\url{https://github.com/Aniket2002/ifrs16-lbo-engine}.
Exact revision and input hashes accompany this paper.
\end{thebibliography}

\newpage\appendix
\section{Mathematical Proofs}\label{app:proofs}
\subsection{Exact ratio perturbation}
For nonzero $y,\widehat y$,
\begin{equation}
\frac{x}{y}-\frac{\widehat x}{\widehat y}
=\frac{x-\widehat x}{y}
 +\frac{\widehat x(\widehat y-y)}{y\widehat y}.
\label{eq:ratioidentity}
\end{equation}
If both denominator magnitudes are at least $b>0$, taking absolute values gives
$|x/y-\widehat x/\widehat y|\le|x-\widehat x|/b+
|\widehat x|\,|\widehat y-y|/b^2$. A floor on only the unperturbed denominator
does not suffice. For example, $x=\widehat x=1$, $y=1$,
$\widehat y=1/2$ gives difference 1, exceeding the erroneous first-order bound
$|\widehat y-y|/y^2=1/2$. Thus the historical inequality cannot be treated as a
finite-perturbation theorem.

\begin{proof}[Proof of Proposition~\ref{prop:screening}]
Apply \eqref{eq:ratioidentity} with $(x,y)=(E,I)$ and then $(N,E)$.
Use the respective common denominator floors $i$ and $e$ and the stated analytic
numerator bounds. The triangle inequality gives
$|N^a-N|\le|D^a-D|+|L^a-L|+|C^a-C|$ for matched definitions.
With nonnegative identical rates and timing,
$|I^a-I|\le r_d|D^a-D|+r_L|L^a-L|$. Substitution yields the result.
\end{proof}
\begin{proof}[Proof of Proposition~\ref{prop:frontier}]
Every $x$ satisfying $\Lev_t^a(x)+\epsilon_{\rm lev}(t)\le c_1^{\rm lev}$
also satisfies the inequality with $c_2^{\rm lev}$. All other constraints and the
objective are unchanged, so set inclusion and monotonicity of the supremum follow.
No continuity or strict improvement is implied.
\end{proof}
\begin{proof}[Proof of Proposition~\ref{prop:certification}]
At each date, $\ICR_t\ge\ICR_t^a-\epsilon_{\rm icr}(t)>c^{\rm icr}$
and $\Lev_t\le\Lev_t^a+\epsilon_{\rm lev}(t)<c^{\rm lev}$.
These are the only events certified by the premise.
\end{proof}

\section{Theoretical Assumptions}\label{app:assumptions}
\begin{assumption}[Matched definitions]\label{ass:growth}
The two theoretical paths have the same testing dates, accounting definitions,
units, recurrence coefficients and interest timing. Any differing cash flows,
initial values or inputs have explicitly established error budgets.
\end{assumption}
\begin{assumption}[Positive denominator floors]\label{ass:margins}
Both EBITDA paths exceed a known $e>0$ and both applicable interest paths exceed
a known $i>0$ at each test. If either floor fails, no ratio certification is made.
\end{assumption}
\begin{assumption}[Finite perturbations]\label{ass:sweep}
Component perturbations are bounded in absolute value; analytic numerator upper
bounds are finite. Cash error enters net debt. Debt, lease and rent conventions
are fixed for each ratio comparison.
\end{assumption}
\begin{assumption}[Reduced-form recurrences]\label{ass:rates}
For the recurrence lemma below, both paths use the same debt coefficient
$a=1+r_d\ge0$, sweep $s\ge0$, and the same nonnegative projection. Lease paths
use a common $q\in[0,1]$ and explicitly bounded addition differences.
\end{assumption}
These assumptions replace incompatible historical growth, CPI and quarterly
testing assertions. They do not claim that the full benchmark meets the
matched-model conditions.

\newpage
\section{Explicit Error Bound Derivations}\label{app:error_bounds}
\begin{lemma}[Matched recurrence stability]\label{lem:error_bounds}
Consider the matched recurrences
\[
D_{t+1}=\max(0,aD_t-sF_t),\qquad
\widehat D_{t+1}=\max(0,a\widehat D_t-s\widehat F_t),
\]
with
$|F_t-\widehat F_t|\le\delta F_t$. Then
\begin{equation}
|D_t-\widehat D_t|\le a^t\epsilon_D(0)
 +s\sum_{k=0}^{t-1}a^{t-1-k}\delta F_k.\label{eq:debtbudget}
\end{equation}
For lease recurrences $L_{t+1}=\max(0,(1-q)L_t+A_t^L)$ with identical $q$,
\begin{equation}
|L_t-\widehat L_t|\le(1-q)^t\epsilon_L(0)
 +\sum_{k=0}^{t-1}(1-q)^{t-1-k}\delta A_k^L.
\label{eq:leasebudget}
\end{equation}
\end{lemma}
\begin{proof}
The map $x\mapsto\max(0,x)$ is 1-Lipschitz. Hence the debt difference at $t+1$
is at most $a$ times the difference at $t$ plus $s\delta F_t$.
Iteration proves \eqref{eq:debtbudget}; replacing $a$ with $1-q$ and the forcing
error with $\delta A_t^L$ proves \eqref{eq:leasebudget}.
\end{proof}
These formulas propagate known forcing errors; they do not establish those
errors from an arbitrary FCF conversion assumption. If a forcing bound has the
form $\delta F_k\le b(1+g)^k$, its sum is
\begin{equation}
\sum_{k=0}^{t-1}a^{t-1-k}(1+g)^k=
\begin{cases}\displaystyle\frac{a^t-(1+g)^t}{a-(1+g)},&a\ne1+g,\\
t a^{t-1},&a=1+g.
\end{cases}
\end{equation}
For a separately stipulated CPI-indexed payment model, preserve the conventional
algebraic schedule $L_{t+1}=(1+r_L)L_t-P_t$, with
$P_t=P_0(1+\pi)^t$. It is not the current run-off implementation. For
$-1<\pi_{\min}\le\pi,\widehat\pi\le\pi_{\max}$, the mean-value theorem yields
\begin{equation}
|P_t(\pi)-P_t(\widehat\pi)|\le
P_0t(1+\pi_{\max})^{t-1}|\pi-\widehat\pi|\quad(t\ge1),
\end{equation}
and payment differences can be propagated as in \eqref{eq:debtbudget} with
coefficient $1+r_L$. This avoids the historical expression with denominator
$1-(1+\pi_{\min})^{-T}$, which is negative for $\pi_{\min}=-0.02$ and $T>0$.
No numeric CPI bound is transferred to the full simulation.

Finally, for constant growth $g,\widehat g\in[g_-,g_+]$ with $g_->-1$ and the
same starting EBITDA,
\begin{equation}
|E_0(1+g)^t-E_0(1+\widehat g)^t|
\le |E_0|t(1+g_+)^{t-1}|g-\widehat g|.
\end{equation}
This finite-perturbation bound follows from the mean-value theorem. A capex
conversion parameter changes FCF, not this EBITDA growth identity. None of these
statements supplies a numerical certification rate for the empirical experiment.

\section{Benchmark Generator Details}\label{app:benchmark}
The archive contains the complete input assumptions for both model families,
scenario-level outcomes and the 1,000 annual simulation and analytic observations.
The figure generator reconstructs the seed-42 draws and checks the analytic
scores, simulated maximum leverage and minimum coverage against the archived
benchmark before plotting. Figures and tables use the current numerical output;
the original figures' filenames are retained to make replacement traceable.

\newpage
\section{Computational Environment and Reproducibility}\label{app:computation}
\begingroup\small\raggedright
\paperinput{provenance.tex}
\endgroup
The source revision identifies the exact committed manuscript and generation
code. A subsequent artifact-only commit stores the generated provenance, figures,
benchmark and PDF: a commit cannot embed its own final hash in its own contents.
The artifact manifest records actual SHA-256 checksums, including the environment
export. No release tag, DOI deposit, historical environment hash or fixed timing
claim is asserted. The repository license is MIT.

From the repository root, with dependencies installed, run:
\begin{verbatim}
python -m ruff check .
python -m ruff format --check .
python -m pytest -q --cov=src/lbo --cov-fail-under=52
python -m analysis.run_benchmark --seed 42
python -m analysis.scripts.generate_paper_v2_figures --run-benchmark
latexmk -pdf -outdir=paper paper/ifrs16_lbo_ssrn_v2.tex
\end{verbatim}
The figure command archives a fresh full report in
\path{results/paper_v2/benchmark_seed42.json}; the standalone benchmark otherwise
overwrites its output directory. When latexmk is absent, use repeated pdflatex
passes or the documented Tectonic fallback in \path{paper/REPRODUCE_V2.md}.
Financial quantities are reproducible for fixed code, inputs, seed and dependencies;
wall-clock measurements are environment-dependent. Local validation and PDF
inspection records are archived separately from the scientific claims.
\end{document}
